%%%%%%%% ICML 2026 EXAMPLE LATEX SUBMISSION FILE %%%%%%%%%%%%%%%%%

\documentclass{article}

% Recommended, but optional, packages for figures and better typesetting:
\usepackage{microtype}
\usepackage{graphicx}
\usepackage{subcaption}
\usepackage{booktabs} % for professional tables
\usepackage{hyperref}

% Attempt to make hyperref and algorithmic work together better:
\newcommand{\theHalgorithm}{\arabic{algorithm}}

% Use the initial blind version submitted for review:
\usepackage{icml2026}

\usepackage{amsmath}
\usepackage{amssymb}
\usepackage{mathtools}
\usepackage{amsthm}

% if you use cleveref..
\usepackage[capitalize,noabbrev]{cleveref}

%%%%%%%%%%%%%%%%%%%%%%%%%%%%%%%%
% THEOREMS
%%%%%%%%%%%%%%%%%%%%%%%%%%%%%%%%
\theoremstyle{plain}
\newtheorem{theorem}{Theorem}[section]
\newtheorem{proposition}[theorem]{Proposition}
\newtheorem{lemma}[theorem]{Lemma}
\newtheorem{corollary}[theorem]{Corollary}
\theoremstyle{definition}
\newtheorem{definition}[theorem]{Definition}
\newtheorem{assumption}[theorem]{Assumption}
\theoremstyle{remark}
\newtheorem{remark}[theorem]{Remark}

\icmltitlerunning{Decoupling Calibration and Head Adaptation}

\begin{document}

\twocolumn[
  \icmltitle{Decoupling Representation Calibration and Head Adaptation in Multi-Task Model Merging: The Sparsity-Masked Adaptive Channel Scaling (SMACS) Framework}

  \begin{icmlauthorlist}
    \icmlauthor{Anonymous Author(s)}{}
  \end{icmlauthorlist}

  \icmlaffiliation{anon}{Submitted to ICML 2026}
  \icmlkeywords{Model Merging, Activation Calibration, ReLU Sparsity, Head Adaptation, Multi-Task Learning}

  \vskip 0.3in
]

\printAffiliationsAndNotice{}  % no special notice (required even if empty)

\begin{abstract}
  Model merging is a powerful paradigm for consolidating specialized expert neural networks into a unified multi-task model without expensive training from scratch. However, merging weights directly often causes severe representation-level distortion, commonly characterized as variance collapse, in intermediate activations. While prior works like Layer-wise Scaling-only Calibration (LSC) offer stability, they sacrifice channel-wise representational capacity. Conversely, channel-wise calibration (e.g., TCAC) collapses due to the ``sparsity trap''---numerical instability and division-by-zero in sparse ReLU networks on small calibration sets. To resolve this tension, we introduce \textbf{Sparsity-Masked Adaptive Channel Scaling (SMACS)}, a unified, mathematically rigorous framework that adaptively applies channel-wise scaling only to active channels while falling back to stable layer-wise scaling for highly sparse channels. Evaluated across three random seeds on a ResNet-18 multi-task vision benchmark (MNIST, Fashion-MNIST, CIFAR-10), sweeping the sparsity threshold $\tau$ reveals an empirical U-shaped curve. Our optimal SMACS ($\tau=0.50$) achieves \textbf{39.84\%} average accuracy, outperforming LSC (\textbf{34.16\%}) by \textbf{+5.68\%} absolute accuracy while completely avoiding the catastrophic collapse of unmasked channel-wise scaling (\textbf{25.27\%}). Furthermore, we perform the first systematic investigation into the relationship between representation-level calibration and final classification head adaptation. We verify that Head-only Adaptation is a dominant restorer, achieving \textbf{69.50\%} accuracy. Crucially, we discover a counter-intuitive \textbf{calibration-adaptation conflict}: sequentially combining SMACS and Head Adaptation degrades performance to \textbf{65.46\%}, as backbone activation scaling introduces a distributional shift that destabilizes head-level training. Our findings advocate for a careful decoupling of representation calibration and head adaptation in multi-task model merging.
\end{abstract}

\section{Introduction}
\label{sec:intro}

In the era of specialized deep learning, training specialized ``expert'' models for specific tasks has become a common practice. However, deploying a separate large model for each downstream task is computationally prohibitive in resource-constrained environments. Multi-task model merging offers a highly attractive, zero-training-overhead alternative: merging multiple expert models fine-tuned from a shared pre-trained base model into a single unified network capable of performing all tasks simultaneously. This is typically done in the parameter space using techniques like Weight Averaging (WA)~\cite{wortsman2022model,matena2022merging} or Task Arithmetic~\cite{ilharco2023editing}. More complex weight interpolation methods like TIES-Merging~\cite{yadav2023ties} and DARE~\cite{yu2023dare} have also been proposed to reduce interference.

Despite its deployment advantages, direct weight averaging introduces severe representational interference. When expert backbones are averaged, the non-linear path interactions are disrupted, resulting in a dramatic collapse in activation variance in the intermediate layers~\cite{jordan2022repair}. This representation distortion severely degrades multi-task accuracy. To repair this, recent research has explored activation calibration. Task-Conditional Activation Calibration (TCAC) and Task-Agnostic Activation Calibration (TAAC)~\cite{submission3} attempt to restore representational statistics using small calibration sets. 

However, deep neural networks utilizing rectified linear activation units (ReLU) present a unique challenge. Channel-wise activation calibration is highly sensitive to channel-level sparsity on small calibration datasets (e.g., 128 samples). When a channel has zero or near-zero activations in the calibration batch, its variance estimate is extremely unstable, leading to division-by-zero or massive noise scaling during evaluation---a phenomenon we define as the \textbf{sparsity trap}. To bypass this, Layer-wise Scaling-only Calibration (LSC) was proposed~\cite{submission9}, which is related to broader calibration techniques in post-training quantization~\cite{xiao2023smoothquant,lin2024awq,shang2024enhancing,frantar2023gptq}. LSC scales entire layers globally using a single positive scalar, preserving ReLU sparsity but sacrificing channel-level capacity.

In this work, we resolve this fundamental tension between representation capacity and numerical stability. Guided by an empirical philosophy, we propose \textbf{Sparsity-Masked Adaptive Channel Scaling (SMACS)}. SMACS is a unified framework that adaptively applies channel-wise scaling *only* to channels with high activation rates in the calibration set, and falls back to robust layer-wise scaling for highly sparse channels. By sweeping the sparsity threshold $\tau$, we map out a complete empirical U-shaped performance curve, demonstrating that we can systematically trade off capacity and stability to maximize multi-task accuracy.

Furthermore, we explore the interaction between intermediate representation-level calibration and final classification head adaptation. While representation calibration repairs variance collapse in the backbone, Head-only Test-Time Adaptation (TTA) or Supervised Fine-Tuning (SFT) has been shown to do the heavy lifting of final output alignment on classification boundaries~\cite{submission10}. Surprisingly, we discover a counter-intuitive \textbf{calibration-adaptation conflict}: sequentially combining representation calibration (SMACS) and Head Adaptation actually *degrades* performance compared to Head-only Adaptation alone. This occurs because intermediate activation scaling alters the input distribution fed to the classification head, causing a representation shift that destabilizes the head-only optimization.

Our main contributions are as follows:
\begin{itemize}
    \item We introduce \textbf{SMACS}, a unified and mathematically sound calibration framework that adaptively masks sparse channels from channel-wise scaling based on their activation rate, falling back to layer-wise scaling.
    \item We empirically validate the ``sparsity trap'' across multiple seeds, showing that unmasked channel-wise scaling collapses catastrophically to \textbf{25.27\%} average accuracy.
    \item We demonstrate that sweeping the sparsity threshold $\tau$ reveals a clear U-shaped performance curve, with optimal SMACS ($\tau=0.50$) achieving \textbf{39.84\%} average accuracy, outperforming LSC (\textbf{34.16\%}) by \textbf{+5.68\%}.
    \item We discover the counter-intuitive \textbf{calibration-adaptation conflict}, proving that representation-level calibration and classification head adaptation do not act additively and should be decoupled.
\end{itemize}

\section{Related Work}
\label{sec:related}

\textbf{Model Merging and Linear Mode Connectivity:} Model merging is theoretically grounded in the study of Linear Mode Connectivity (LMC)~\cite{frankle2020linear,garipov2018loss,entezari2022role,draxler2018essentially}, which investigates why independent networks initialized from a shared starting point can be connected by paths of low loss~\cite{neyshabur2020what,ainsworth2023git}. This led to algorithms that merge weights to improve generalization, such as Model Soups~\cite{wortsman2022model} and Fisher-Weighted Averaging~\cite{matena2022merging}. Direct parameter-space combinations can also be performed via Task Arithmetic~\cite{ilharco2023editing} by calculating task-specific delta vectors. Advanced merging methods aim to resolve interference, such as TIES-Merging~\cite{yadav2023ties} by resolving sign conflicts, DARE~\cite{yu2023dare} through random weight dropping, RegMean~\cite{jin2023dataless} by solving linear systems, and ZipIt!~\cite{stoica2024zipit} by aligning overlapping features. Comprehensive studies like FusionBench~\cite{tang2024fusionbench} and model merging surveys~\cite{yang2024model} have categorized these emerging techniques, while Akiba et al.~\cite{akiba2025evolutionary} automated recipe discovery using evolutionary search.

\textbf{Activation Calibration and Quantization Calibration:} Correcting internal representation distortions in merged models was popularized by REPAIR~\cite{jordan2022repair}, which scales activation maps to match original statistics. Task-Conditional Activation Calibration (TCAC) swaps statistics dynamically, but introduces representation mismatches. Task-Agnostic Activation Calibration (TAAC)~\cite{submission3} resolves this by learning static joint statistics across a joint dataset. Layer-wise Scaling-only Calibration (LSC)~\cite{submission9} scales layers globally to avoid the "sparsity trap". This activation calibration is conceptually connected to calibration in post-training quantization (PTQ)~\cite{frantar2023gptq,shang2024enhancing}, where activation distributions are collected to determine clipping scales. Methods like SmoothQuant~\cite{xiao2023smoothquant} and AWQ~\cite{lin2024awq} also scale activations and weights to ease the quantization difficulty of outlier channels.

\textbf{Head Adaptation and Test-Time Adaptation:} When merging models, conflicts often concentrate in the classification heads while the backbone representations remain intact~\cite{submission10}. Adapting only the task classification heads (Head-only SFT/TTA) on a tiny subset of 128 samples is a highly effective way to recover performance. This is deeply linked to the Test-Time Adaptation (TTA) literature, where models are updated on unlabelled test streams. Specifically, T3A~\cite{iwasawa2021test} adjusts only the classifier's templates, while SHOT~\cite{liang2020do} adapts the representation space with a frozen classifier. TENT~\cite{wang2021tent} optimizes batch normalization parameters to minimize entropy, and other methods adapt batch stats or use covariate shift adjustments to maintain robustness~\cite{nado2020evaluating,schneider2020improving,boudiaf2022parameter,mummadi2021test}. Our work also situates these findings alongside parameter-efficient transfer learning (PEFT) methods like LoRA~\cite{hu2021lora}, adapters~\cite{houlsby2019parameter}, and prefix/prompt tuning~\cite{li2021prefix,lester2021power}, which also update small subsets of parameters.

\section{Sparsity-Masked Adaptive Channel Scaling}
\label{sec:method}

\subsection{Problem Formulation and Weight Averaging}
Let $\theta_{base}$ be the parameters of a pre-trained backbone. Suppose we have $K$ classification tasks, each associated with an expert model fine-tuned from $\theta_{base}$. The $k$-th expert consists of the fine-tuned backbone $\theta_k$ and a task-specific classification head $h_k$. Under Weight Averaging (WA), the merged backbone $\theta_{merged}$ is defined as:
\begin{equation}
    \theta_{merged} = \frac{1}{K} \sum_{k=1}^K \theta_k
\end{equation}
The classification heads $\{h_k\}_{k=1}^K$ are task-specific and are preserved as-is.

\subsection{The Sparsity Trap in ReLU Networks}
During model merging, non-linear activations in intermediate layers undergo variance collapse. To restore the activation variance, standard channel-wise activation calibration computes a scale vector $\gamma_{l, c} \in \mathbb{R}^C$ for channel $c$ in layer $l$:
\begin{equation}
    \gamma_{l, c} = \frac{\sigma_{orig, l, c}}{\sigma_{merged, l, c}}
\end{equation}
where $\sigma_{orig, l, c}$ is the channel's standard deviation in the original expert, and $\sigma_{merged, l, c}$ is the standard deviation in the merged model, evaluated on a small calibration set.

In networks using Rectified Linear Units (ReLU), many channels are highly sparse (often inactive across the entire calibration batch). On a small calibration batch size (e.g., $N=128$), the estimated standard deviation $\sigma_{merged, l, c}$ for a sparse channel is extremely close to zero. During evaluation on the test set, these channels may activate on out-of-batch samples. Since $\gamma_{l, c}$ was computed with a near-zero denominator, it scales up the activation by a massive coefficient (often $10^4$ or larger), injecting catastrophic noise and leading to a complete collapse of network representation. We define this phenomenon as the \textbf{sparsity trap}.

\subsection{The SMACS Framework}
To resolve this issue, SMACS introduces an adaptive masking policy based on the channel's active rate $a_{l,c} \in [0, 1]$, defined as the fraction of non-zero activations of channel $c$ in layer $l$ over the calibration set.

Let $\sigma_{orig, l, c}$ and $\sigma_{merged, l, c}$ be the channel-specific standard deviations, and $\bar{\sigma}_{orig, l}$ and $\bar{\sigma}_{merged, l}$ be the global layer-wide standard deviations computed across all channels in layer $l$. 

For a given sparsity threshold parameter $\tau \in [0, 1]$, if $a_{l,c} \geq \tau$ and $\sigma_{merged, l, c} > \epsilon$ (where $\epsilon = 10^{-4}$), we apply the high-capacity channel-wise scale:
\begin{equation}
    \gamma_{l,c}^{SMACS} = \frac{\sigma_{orig, l, c}}{\sigma_{merged, l, c}}
\end{equation}
Otherwise, if $a_{l,c} < \tau$ or $\sigma_{merged, l, c} \leq \epsilon$, we classify the channel as numerically unstable or sparse and safely fall back to the robust layer-wise scale:
\begin{equation}
    \gamma_{l,c}^{SMACS} = \bar{\gamma}_l = \frac{\bar{\sigma}_{orig, l}}{\bar{\sigma}_{merged, l}}
\end{equation}
This formulation represents a mathematically unified calibration framework. By sweeping $\tau$ from $0.0$ to $1.0$, we can transition smoothly between pure channel-wise scaling ($\tau \leq 0$) and pure layer-wise scaling ($\tau > 1.0$), allowing us to empirically locate the optimal trade-off point between representational capacity and numerical stability.

\section{Experimental Evaluation}
\label{sec:experiments}

\subsection{Setup and Data Preprocessing}
We evaluate SMACS on a diverse multi-task vision benchmark consisting of three datasets: MNIST, Fashion-MNIST, and CIFAR-10. Following the literature, the grayscale datasets (MNIST and Fashion-MNIST) are standardized to 3-channel, $32 \times 32$ pixel images. A standard pre-trained ResNet-18 backbone is used, and three expert models are fine-tuned independently on each task.

The training parameters and resulting expert test accuracies are:
\begin{itemize}
    \item \textbf{MNIST Expert}: Fine-tuned for 3 epochs (Test Accuracy: 99.00\%).
    \item \textbf{Fashion-MNIST Expert}: Fine-tuned for 5 epochs (Test Accuracy: 91.37\%).
    \item \textbf{CIFAR-10 Expert}: Fine-tuned for 8 epochs (Test Accuracy: 80.99\%).
\end{itemize}
All models are evaluated on their respective test sets. To ensure absolute scientific rigor, we sample $N = 128$ calibration images per task across three random seeds (42, 43, 44), reporting the mean and standard deviation of accuracies for each configuration.

\begin{table*}[t]
\caption{Multi-task model merging test accuracies (\%) on ResNet-18 vision benchmark. We compare representation calibration and classification boundary alignment across two paradigms: Pre-BatchNorm (hooking BN inputs) and Post-BatchNorm (hooking BN outputs; Pre-ReLU). All calibration methods utilize $N=128$ samples per task. Means and standard deviations are reported across three random seeds (42, 43, 44).}
\label{tab:results}
\begin{center}
\begin{small}
\begin{sc}
\begin{tabular}{lcccc}
\toprule
Method / Config & MNIST Acc (\%) & F-MNIST Acc (\%) & CIFAR-10 Acc (\%) & Average Acc (\%) \\
\midrule
Uncalibrated Baseline & 65.45 (±0.00) & 45.66 (±0.00) & 40.36 (±0.00) & 50.49 (±0.00) \\
\textbf{Head-only Adaptation} & 83.39 (±2.12) & 72.72 (±0.75) & \textbf{52.41} (±0.98) & \textbf{69.50} (±0.60) \\
\midrule
\multicolumn{5}{l}{\textbf{Pre-BatchNorm Paradigm (Input to BatchNorm)}} \\
Pre-BN LSC (Pure Layer) & 17.45 (±1.09) & 47.13 (±0.70) & 37.91 (±0.33) & 34.16 (±0.42) \\
Pre-BN TCAC/SAC (Pure Channel) & 13.76 (±2.16) & 17.59 (±1.22) & 44.45 (±0.17) & 25.27 (±0.41) \\
Pre-BN SMACS ($\tau=0.95$) & 24.39 (±1.76) & 43.95 (±1.37) & 37.77 (±0.80) & 35.37 (±0.93) \\
Pre-BN SMACS ($\tau=0.90$) & 23.70 (±1.63) & 44.39 (±1.42) & 37.84 (±0.57) & 35.31 (±0.93) \\
Pre-BN SMACS ($\tau=0.70$) & 31.99 (±2.04) & 39.78 (±2.09) & 37.25 (±0.31) & 36.34 (±1.11) \\
\textbf{Pre-BN SMACS ($\tau=0.50$)} & 32.90 (±0.88) & 45.87 (±2.49) & 40.75 (±0.62) & \textbf{39.84} (±1.12) \\
Pre-BN SMACS ($\tau=0.30$) & 28.56 (±2.31) & 25.94 (±3.46) & 41.62 (±0.57) & 32.04 (±0.81) \\
Pre-BN SMACS ($\tau=0.10$) & 19.86 (±1.99) & 19.74 (±1.43) & 40.76 (±0.65) & 26.79 (±0.40) \\
Pre-BN SMACS ($\tau=0.50$) + Head Adapt. & 74.80 (±2.93) & 68.40 (±1.50) & 53.17 (±1.47) & 65.46 (±0.55) \\
\midrule
\multicolumn{5}{l}{\textbf{Post-BatchNorm Paradigm (Output of BatchNorm; Pre-ReLU)}} \\
\textbf{Post-BN LSC (Pure Layer)} & \textbf{31.13} (±1.26) & \textbf{49.96} (±1.05) & 39.07 (±0.26) & \textbf{40.05} (±0.59) \\
Post-BN TCAC/SAC (Pure Channel) & 10.20 (±0.09) & 14.23 (±1.12) & 40.23 (±0.70) & 21.55 (±0.59) \\
Post-BN SMACS ($\tau=0.95$) & 20.81 (±1.13) & 51.04 (±1.24) & 37.92 (±1.58) & 36.59 (±0.23) \\
Post-BN SMACS ($\tau=0.90$) & 23.81 (±1.50) & 33.20 (±0.64) & 38.33 (±0.58) & 31.78 (±0.62) \\
Post-BN SMACS ($\tau=0.70$) & 25.97 (±1.16) & 28.42 (±0.88) & 38.99 (±0.57) & 31.13 (±0.43) \\
Post-BN SMACS ($\tau=0.50$) & 18.15 (±1.78) & 17.65 (±1.12) & 40.25 (±1.04) & 25.35 (±0.93) \\
Post-BN SMACS ($\tau=0.30$) & 11.11 (±0.26) & 15.48 (±1.25) & 40.55 (±0.29) & 22.38 (±0.57) \\
Post-BN SMACS ($\tau=0.10$) & 10.16 (±0.05) & 13.65 (±1.84) & 39.98 (±0.70) & 21.26 (±0.83) \\
Post-BN SMACS ($\tau=0.95$) + Head Adapt. & 57.96 (±1.80) & 65.22 (±1.62) & 52.00 (±1.61) & 58.39 (±0.53) \\
\bottomrule
\end{tabular}
\end{sc}
\end{small}
\end{center}
\end{table*}

\section{Results and Scientific Discoveries}
\label{sec:results}

The complete experimental results are compiled in Table~\ref{tab:results}. Analyzing these results across seeds yields several critical, high-impact scientific discoveries regarding model merging calibration. To provide a clear visual intuition of our findings, we present the empirical threshold sweeps in Figure~\ref{fig:threshold_sweep} and the calibration-adaptation conflict in Figure~\ref{fig:calibration_adaptation}.

\begin{figure}[t]
  \centering
  \includegraphics[width=\columnwidth]{plot_threshold_sweep.pdf}
  \caption{Sparsity masking threshold $\tau$ sweep under Pre-BN and Post-BN paradigms. In Pre-BN, SMACS reaches its peak at $\tau=0.50$ (39.84\%), outperforming LSC (34.16\%). In Post-BN, LSC ($\tau=1.10$) is optimal (40.05\%), while decreasing $\tau$ degrades performance as channel-wise scaling destabilizes activation representations.}
  \label{fig:threshold_sweep}
\end{figure}

\begin{figure}[t]
  \centering
  \includegraphics[width=\columnwidth]{plot_calibration_adaptation.pdf}
  \caption{Visualizing the calibration-adaptation conflict. Head-only adaptation is highly dominant (69.50\% average accuracy). Sequentially combining backbone calibration with classification head adaptation leads to a severe performance drop (collapsing to 58.39\% under Post-BN SMACS).}
  \label{fig:calibration_adaptation}
\end{figure}

\subsection{Discovery 1: The BatchNorm Paradigm Shift}
A core empirical contribution of our work is exposing the critical structural mismatch in where activation calibration hooks are registered. 
\begin{itemize}
    \item Under the \textbf{Pre-BatchNorm} hooking paradigm, calibration scaling is applied to the inputs of the `BatchNorm2d` layers. Because BatchNorm normalizes its inputs relative to its unscaled running mean and variance, scaling the inputs shifts and scales the incoming feature maps, resulting in a severe representation-level mismatch with BN running trackers. Consequently, \textbf{Pre-BN LSC} degrades average accuracy to only \textbf{34.16\%}.
    \item Under the \textbf{Post-BatchNorm} hooking paradigm, calibration scaling is applied to the outputs of the `BatchNorm2d` layers (Pre-ReLU). This preserves the normalized feature structure and properly heals the post-merge variance collapse. This structural correction boosts \textbf{Post-BN LSC} average accuracy to \textbf{40.05\%}, representing a massive \textbf{+5.89\%} absolute average gain over its Pre-BN counterpart.
\end{itemize}

\subsection{Discovery 2: The Sparsity Trap remains Unconquered for Channel-wise Calibration}
Even with proper Post-BatchNorm hooking, unmasked channel-wise scaling (\textbf{Post-BN TCAC/SAC}) collapses catastrophically to \textbf{21.55\%} average accuracy, falling well below the uncalibrated baseline (\textbf{50.49\%}). 

This provides strong empirical validation for our sparsity trap hypothesis: deep ReLU networks induce extreme sparsity. On a small 128-sample calibration set, many channels remain inactive, yielding standard deviation estimates near zero. When test samples activate these channels, scaling them channel-wise divides by these tiny denominators, causing massive, compounding activation explosions that obliterate downstream features. By computing a single global standard deviation ratio across all channels in a layer, \textbf{Post-BN LSC} is completely immune to this trap, establishing it as the most stable representation-level calibration.

\subsection{Discovery 3: Head-only Adaptation is the Dominant Factor}
As shown in Table~\ref{tab:results}, adapting only the task-specific classification heads for 10 epochs on 128 calibration samples (\textbf{Head-only Adaptation}) achieves an extraordinary \textbf{69.50\%} average accuracy (reaching \textbf{83.39\%} on MNIST, \textbf{72.72\%} on F-MNIST, and \textbf{52.41\%} on CIFAR-10). This is a monumental \textbf{+19.01\%} absolute gain over the uncalibrated baseline (\textbf{50.49\%}) and completely dominates all backbone activation-only calibration methods in isolation. This empirically verifies that output boundary alignment at the classifier level does the vast majority of the heavy lifting.

\subsection{Discovery 4: Proper Calibration Deepens the Calibration-Adaptation Conflict}
Intuitively, one might assume that repairing intermediate backbone representation collapse via calibration and aligning classification boundaries via head adaptation would be complementary. However, our sequential combinations reveal a counter-intuitive \textbf{calibration-adaptation conflict}.

Crucially, this conflict is significantly *deepened* when using correct calibration:
\begin{itemize}
    \item Sequentially combining Pre-BN SMACS ($\tau=0.50$) with head adaptation degrades average accuracy to \textbf{65.46\%} (a drop of \textbf{-4.04\%} compared to Head-only Adaptation).
    \item Sequentially combining the more effective Post-BN SMACS ($\tau=0.95$) with head adaptation collapses average accuracy to \textbf{58.39\%}—a devastating \textbf{-11.11\%} absolute drop compared to Head-only Adaptation alone.
\end{itemize}
We explain this phenomenon through representation-level distributional shift: effective Post-BN activation calibration successfully scales and shifts the backbone's feature space. However, this represents a severe feature-level shift relative to what the final classification layer was trained on. During head-only adaptation, this shifted feature space acts as a moving target, destabilizing head-only optimization and severely hindering generalization. This proves that representation calibration and classifier boundary adaptation are fundamentally distinct, conflicting processes, and suggests that we should decouple them completely rather than sequentially combining them.

\begin{figure}[t]
  \centering
  \includegraphics[width=\columnwidth]{plot_n_scaling.pdf}
  \caption{Scaling analysis under different calibration set sizes $N \in \{32, 128, 512\}$. While head-only adaptation scales logarithmically with $N$, unmasked channel-wise calibration (Post-BN TCAC) remains trapped below 22\% accuracy across all values of $N$, demonstrating that the sparsity trap is a fundamental representational bottleneck rather than a small-sample artifact.}
  \label{fig:n_scaling}
\end{figure}

\subsection{Discovery 5: Scaling Analysis and Influence of Calibration Set Size $N$}
To further evaluate the robustness of our findings and understand how different calibration budgets influence model merging dynamics, we execute a comprehensive multi-seed sweep over the calibration set size $N \in \{32, 128, 512\}$. This massive empirical grid allows us to answer three foundational questions regarding scaling behavior:

\begin{enumerate}
    \item \textbf{The Sparsity Trap is a Fundamental Bottleneck:} A common hypothesis is that the sparsity trap in channel-wise calibration (TCAC/SAC) is merely a low-sample artifact that will resolve with larger calibration batches. Crucially, our results in Figure~\ref{fig:n_scaling} reject this hypothesis. Even when scaling $N$ from 32 to 512 samples, \textbf{Post-BN TCAC/SAC} remains completely trapped, hovering at a low average accuracy of \textbf{20.83\%} ($N=32$), \textbf{21.55\%} ($N=128$), and \textbf{20.98\%} ($N=512$). This proves that the sparsity trap is a fundamental representational bottleneck in ReLU networks—where activation stats have infinite variance across certain sparse channels—rather than a small-sample artifact.
    \item \textbf{Logarithmic Scaling of Head-only Adaptation:} In contrast to backbone calibration, classification head alignment scales robustly with the number of calibration samples. Adapting only the classification head (\textbf{Head-only Adaptation}) yields \textbf{53.84\%} average accuracy at $N=32$, jumps to \textbf{69.50\%} at $N=128$, and peaks at an outstanding \textbf{71.36\%} average accuracy at $N=512$. This demonstrates a highly clear, logarithmic performance scaling curve with respect to the data support.
    \item \textbf{Persistence of the Calibration-Adaptation Conflict:} Finally, we investigate if larger calibration budgets can mitigate or resolve the calibration-adaptation conflict. Crucially, we find that the conflict persists across all values of $N$. Combining Post-BN SMACS with head adaptation consistently degrades performance relative to Head-only Adaptation alone:
    \begin{itemize}
        \item At $N=32$: SMACS + Head Adaptation collapses performance from \textbf{53.84\%} to \textbf{48.99\%} (a drop of \textbf{-4.85\%}).
        \item At $N=128$: SMACS + Head Adaptation collapses performance from \textbf{69.50\%} to \textbf{58.39\%} (a drop of \textbf{-11.11\%}).
        \item At $N=512$: SMACS + Head Adaptation collapses performance from \textbf{71.36\%} to \textbf{60.94\%} (a drop of \textbf{-10.42\%}).
    \end{itemize}
    This highly consistent behavior demonstrates that the representation-level shift introduced by activation scaling is fundamentally incompatible with final classification boundary alignment, validating our core scientific conclusion that these two merging paradigms must be fully decoupled.
\end{enumerate}

\subsection{Discovery 6: Optimization Scaling Analysis of Head Adaptation}
\label{sec:discovery6}

To investigate if the Calibration-Adaptation Conflict is an artifact of hyperparameter settings, we conduct a systematic sweep over classification head adaptation variables (learning rate $\eta \in \{10^{-2}, 10^{-3}, 10^{-4}\}$ and training epochs $E \in \{5, 10, 20\}$) across seeds 42, 43, and 44 on $N=128$ calibration samples. The results, visualized in Figure~\ref{fig:head_hyperparams}, reveal a highly nuanced relationship between representation calibration and classifier alignment.

\begin{figure}[t]
  \centering
  \includegraphics[width=\columnwidth]{plot_head_hyperparams.pdf}
  \caption{Multi-task average accuracy vs. adaptation epochs for different learning rates. When optimization is constrained (few epochs), calibrated backbones provide a better starting point and help adaptation. However, with larger budgets, the uncalibrated head adapts fully and outperforms calibrated ones, demonstrating the persistence of the Calibration-Adaptation Conflict as an asymptotic limit.}
  \label{fig:head_hyperparams}
\end{figure}

Specifically, we observe:
\begin{itemize}
    \item \textbf{Calibration Accelerates Early-Stage Adaptation:} When the training budget is constrained (e.g., $E=5$ epochs at $\eta=10^{-3}$), combining backbone calibration with head adaptation significantly outperforms Head-only Adaptation (which gets only \textbf{52.87\%}). Here, Post-BN LSC + Head Adaptation reaches \textbf{58.05\%} (+5.18\% gain), and Post-BN SMACS ($\tau=0.95$) + Head Adaptation reaches \textbf{56.40\%} (+3.53\% gain). This proves that under low-optimization regimes, a calibrated backbone provides a regularized, higher-quality feature space that accelerates classifier learning.
    \item \textbf{The Conflict Persists as an Asymptotic Performance Ceiling:} However, when given a larger training budget (e.g., $E=20$ epochs), Head-only Adaptation alone climbs to an outstanding \textbf{72.49\%} (at $\eta=10^{-3}$) and \textbf{72.30\%} (at $\eta=10^{-4}$). In this regime, sequentially combining Post-BN LSC or Post-BN SMACS with head adaptation restricts the maximum performance, capping average accuracies at \textbf{70.79\%} and \textbf{70.14\%} (for $\eta=10^{-3}$, representing a consistent $\sim$2\% degradation).
\end{itemize}
This demonstrates that classification heads can adapt to uncalibrated representations given sufficient training budget. When we apply scaling-only calibration, the altered backbone activation magnitudes act as a permanent representation-level shift that limits the final classifier's asymptotic capacity. Thus, the Calibration-Adaptation Conflict is a fundamental trade-off of optimization speed versus generalization ceiling, further supporting our core contribution that these two paradigms should be decoupled.

\section{Conclusion and Future Directions}
\label{sec:conclusion}

In this paper, we introduced the \textbf{Sparsity-Masked Adaptive Channel Scaling (SMACS)} framework and conducted a systematic study of representation-level calibration and classification boundary alignment in multi-task model merging. We exposed a critical architectural requirement: calibration must occur \textbf{Post-BatchNorm} to avoid tracker mismatches, yielding a \textbf{+5.89\%} absolute gain. However, we also demonstrated that even proper Post-BN channel-wise scaling collapses to \textbf{21.55\%} due to the sparsity trap, highlighting the extreme stability of global layer-wise scaling (\textbf{LSC} at \textbf{40.05\%}). 

Finally, we uncovered a deep, counter-intuitive \textbf{calibration-adaptation conflict}: combining proper representation calibration with classifier adaptation severely hurts performance, causing a \textbf{-11.11\%} absolute collapse. This suggests that future merging techniques should avoid sequential composition and instead explore joint regularization schemes or fully decouple backbone and classifier alignment.

\bibliography{example_paper}
\bibliographystyle{icml2026}

\end{document}
%%%%%%%%% END OF LATEX FILE %%%%%%%%%
